\section{AutoGen: framework de programación de conversaciones entre múltiples Agentes}

\subsection{Filosofía central de diseño}

El diseño de AutoGen se basa en una intuición central: \textbf{la conversación (conversation) es el primitivo más básico de interacción entre Agentes}.

\begin{knowledgebox}{Por qué elegir la conversación como primitivo básico}
La conversación es la forma básica en que los seres humanos avanzan en el aprendizaje y demuestran teoremas. El éxito de ChatGPT muestra que las conversaciones de varios turnos se adaptan de manera natural a la mejora iterativa y a la resolución colaborativa. Un diseño centrado en la conversación puede derivar de forma natural todos los demás patrones de interacción entre Agentes.
\end{knowledgebox}

Construir una aplicación de Agentes solo requiere dos pasos:
\begin{enumerate}
    \item \textbf{Definir los Agentes}: cada Agente puede usar un LLM, herramientas o entrada humana como backend
    \item \textbf{Hacer que los Agentes conversen}: orquestar la colaboración mediante patrones de conversación (conversación entre dos, conversación anidada, chat grupal, etc.)
\end{enumerate}

\subsection{La abstracción Conversable Agent}

La abstracción unificada de AutoGen, \textbf{Conversable Agent}, unifica distintos tipos de entidades (LLM, herramientas, seres humanos) como Agentes capaces de conversar:
\begin{itemize}
    \item Se puede elegir un LLM como backend (por ejemplo, GPT-4, Claude)
    \item Se pueden usar herramientas (ejecución de código Python, llamadas a API, etc.)
    \item Se puede conectar entrada humana
    \item Se pueden combinar varios backends
    \item Admite anidamiento: un Agente puede contener en su interior conversaciones de otros Agentes
\end{itemize}

\subsection{Patrones de conversación}

\begin{itemize}
    \item \textbf{Conversación entre dos}: el patrón más simple, ya admite Reflection (un Agente escritor + un Agente revisor que mejoran de forma iterativa)
    \item \textbf{Sequential Chat}: varias conversaciones entre dos encadenadas, como revisiones anidadas desde múltiples perspectivas (revisión de SEO → revisión legal → revisión ética → metarrevisión)
    \item \textbf{Nested Chat}: un Agente inicia subconversaciones en su interior, como cuando el Commander Agent primero conversa con el Writer para obtener el código y después conversa con el Safeguard para verificar su seguridad
    \item \textbf{Group Chat}: varios Agentes con roles distintos eligen dinámicamente al que habla dentro de un grupo, con soporte para condiciones de restricción y lógica de transición de estados
\end{itemize}

\subsection{Ventajas de Multi-Agent}

\begin{importantbox}{¿Por qué múltiples Agentes son mejores que uno solo?}
Los experimentos muestran que descomponer una tarea entre Agentes especializados (como Writer + Safeguard) produce mejores resultados que hacer que un solo Agente atienda todas las instrucciones a la vez: el esquema de múltiples Agentes con GPT-4 supera al esquema de un solo Agente en 20\% en las verificaciones de seguridad, y con GPT-3.5 la diferencia es aún mayor. Cuanto más compleja es la tarea y más débil el modelo, más evidente es la ventaja de múltiples Agentes.
\end{importantbox}

Desde el punto de vista de la programación, la arquitectura de múltiples Agentes también ofrece:
\begin{itemize}
    \item Mejor modularidad y mantenibilidad
    \item Soporte natural para la participación humana (un humano puede tomar el lugar de cualquier rol de Agente)
    \item Facilidad de extensión (reemplazar un solo Agente no afecta al conjunto)
\end{itemize}

\subsection{AutoBuild: construcción automática de equipos de Agentes}

AutoBuild resuelve el problema de ``qué equipo de Agentes diseñar para una tarea específica'':
\begin{enumerate}
    \item El usuario describe el requerimiento de alto nivel
    \item El sistema sugiere automáticamente Agentes con distintos roles
    \item Los Agentes forman un Group Chat para colaborar en la tarea
    \item Adaptive Build además admite construir el equipo de Agentes de forma dinámica y por pasos
\end{enumerate}

\subsection{Resumen de la sección}

AutoGen toma la conversación como primitivo central y ofrece un framework de programación de múltiples Agentes flexible y potente. Su filosofía de diseño es ``definir Agentes + hacer que conversen''; mediante el anidamiento y la composición, es posible construir de manera recursiva sistemas de Agentes de complejidad arbitraria.

